\documentclass[tikz,border=8pt]{standalone}
\usepackage{fontspec}
\usepackage{xeCJK}
\IfFontExistsTF{Noto Sans TC}{\setmainfont{Noto Sans TC}\setCJKmainfont{Noto Sans TC}}{\setmainfont{Microsoft JhengHei}\setCJKmainfont{Microsoft JhengHei}}
\xeCJKsetup{CJKglue={\hskip 0.07em plus 0.02em}, CJKecglue={\hskip 0.10em plus 0.02em}}
\usetikzlibrary{arrows.meta,positioning,calc}
\definecolor{paper}{HTML}{FFFDF8}
\definecolor{navy}{HTML}{0F172A}
\definecolor{muted}{HTML}{334155}
\definecolor{line}{HTML}{CBD5E1}
\definecolor{bad}{HTML}{B91C1C}
\definecolor{badbg}{HTML}{FFE1E1}
\definecolor{good}{HTML}{007A5A}
\definecolor{goodbg}{HTML}{DCFCE7}
\definecolor{blue}{HTML}{1D4ED8}
\definecolor{bluebg}{HTML}{DCEBFF}
\definecolor{gold}{HTML}{B85C00}
\definecolor{goldbg}{HTML}{FFE8B5}
\tikzset{
  card/.style={rounded corners=9pt,line width=1pt,align=center,inner xsep=12pt,inner ysep=7pt,font=\fontsize{8.7}{11.8}\selectfont,text=navy},
  title/.style={font=\bfseries\fontsize{18}{24}\selectfont,text=navy,align=center},
  note/.style={font=\fontsize{10.2}{14}\selectfont,text=muted,align=center}
}
\begin{document}
\begin{tikzpicture}[x=1cm,y=1cm,>=Latex,line cap=round,line join=round]
\path[fill=paper] (0,0) rectangle (16,9.6);
\node[title,text width=14cm] at (8,8.82) {長脈絡不是把所有東西倒回 prompt};
\node[note,text width=13.4cm] at (8,8.02) {我們需要的是可維護的脈絡帳本，而不是越堆越高的流水帳。};

\draw[rounded corners=16pt,draw=bad,line width=1pt,fill=badbg] (.8,1.35) rectangle (6.8,7.25);
\node[font=\bfseries\fontsize{12}{16}\selectfont,text=navy,align=center] at (3.8,6.72) {流水帳式 prompt};
\foreach \y/\t in {5.78/工具回應全文,5.20/模型草稿全文,4.62/錯誤訊息全文,4.04/重複提醒,3.46/過期假設,2.88/上一輪雜訊}{
  \node[card,fill=white,draw=line,text width=4.45cm,minimum height=.42cm] at (3.8,\y) {\t};
}
\node[font=\fontsize{8.9}{12}\selectfont,text=muted,align=center,text width=4.8cm] at (3.8,1.92) {看似記得很多，實際上每一輪都更難判斷。};

\draw[rounded corners=16pt,draw=good,line width=1pt,fill=goodbg] (9.2,1.35) rectangle (15.2,7.25);
\node[font=\bfseries\fontsize{12}{16}\selectfont,text=navy,align=center] at (12.2,6.72) {可維護的脈絡帳本};
\node[card,fill=bluebg,draw=blue,text width=1.65cm,minimum height=1.02cm] at (10.55,5.2) {任務紀錄\\做了什麼};
\node[card,fill=goldbg,draw=gold,text width=1.65cm,minimum height=1.02cm] at (13.85,5.2) {失敗紀錄\\哪裡錯};
\node[card,fill=goodbg,draw=good,text width=1.65cm,minimum height=1.02cm] at (10.55,3.6) {可重用規則\\下次檢查};
\node[card,fill=white,draw=line,text width=1.65cm,minimum height=1.02cm] at (13.85,3.6) {待查問題\\不能裝懂};
\node[font=\fontsize{8.9}{12}\selectfont,text=muted,align=center,text width=4.9cm] at (12.2,2.12) {不是塞更多字，而是留下會再次派上用場的東西。};

\draw[->,draw=navy,line width=1.15pt] (7.15,4.3)--(8.85,4.3);
\node[font=\fontsize{9.5}{12.4}\selectfont,text=muted,align=center,text width=1.65cm] at (8,4.82) {整理};
\node[font=\fontsize{9.5}{12.4}\selectfont,text=muted,align=center,text width=1.65cm] at (8,3.78) {刪減};
\end{tikzpicture}
\end{document}
